% Generated draft: requires mathematical and source review.
\documentclass[11pt]{article}
\usepackage[margin=25mm]{geometry}
\usepackage{fontspec}
\setmainfont{Latin Modern Roman}
\usepackage{amsmath,amssymb,mathtools,unicode-math}
\setmathfont{Latin Modern Math}
\usepackage{enumitem,needspace,xurl,hyperref}
\setlength{\parindent}{0pt}
\setlength{\parskip}{5pt}
\setlength{\emergencystretch}{3em}
\newcommand{\N}{\mathbb{N}}
\newcommand{\Z}{\mathbb{Z}}
\newcommand{\Q}{\mathbb{Q}}
\newcommand{\R}{\mathbb{R}}
\newcommand{\C}{\mathbb{C}}
\newcommand{\F}{\mathbb{F}}
\newcommand{\A}{\mathbb{A}}
\newcommand{\PP}{\mathbb{P}}
\newcommand{\E}{\mathbb{E}}
\newcommand{\eps}{\varepsilon}
\newcommand{\dd}{\,\mathrm{d}}
\begin{document}
\section*{UnsolvedMath ID: 10600041}
\textbf{Source problem code:} catalogue entry
\subsection{Definitions and mathematical statement}
Let \(\mathcal K_{\rm cl}\) and \(\mathcal K_{\rm virt}\) denote classical and virtual knot classes in compatible orientation/framing conventions. A classical diagram determines a virtual diagram. For a quantum invariant \(I\) with value set \(R\), an extension to all virtual knots would assign \(\widetilde I\) on \(\mathcal K_{\rm virt}\), invariant under all required virtual-knot moves, with
\[
 \widetilde I(K)=I(K)\qquad(K\text{ classical}).
\]
Problem-list item 41 asks which quantum invariants admit such an unrestricted extension. It explicitly gives the Kauffman bracket and Jones polynomial as examples. The bracket requires its regular-isotopy/framing convention or normalization, which the record does not state.

The intended extension must be specified more strongly than a bare set-theoretic assignment if it is to classify meaningful quantum constructions: the source gives no permitted local rules, functoriality, algebraic data, or target constraints. Nor does it define the universe of quantum invariants under consideration.

The supplied partial-progress summary says that quantum invariants extend to rotational virtual knots, citing Kauffman--Manturov (2015), and calls that a natural domain. In the usual rotational convention a virtual first-Reidemeister curl move is not imposed, so this is a different equivalence from ordinary unrestricted virtual-knot theory. Consequently that source-reported rotational extension does not, without further compatibility, answer the question for all virtual moves. The summary's examples of bracket and Jones extensions are retained, but a general criterion for ordinary virtual knots is not supplied.
\subsection{Short English statement}
Identify quantum invariants that extend to ordinary virtual knots while respecting every required move and recovering the classical invariant. The record lists bracket and Jones examples, while its progress summary discusses rotational virtual knots, whose move relation differs. The permitted kind of quantum extension remains unspecified.
\subsection{Sources}
\begin{itemize}
\item Local UnsolvedMath registry, permanent ID \(10600041\), exact source code \texttt{catalogue entry}, ``Virtual-knot problem 41.'' No individual proposer or proposed year is supplied.
\item \emph{Unsolved Problems in Virtual Knot Theory and Combinatorial Knot Theory}, revised survey, dedicated problem-list item 41: \url{https://arxiv.org/abs/1409.2823}. Background: 60 top-level items, compound items retained intact; source-TeX extraction, accessed 2026-07-29.
\item Supplied research summary names Kauffman--Manturov (2015) for rotational-virtual extensions but gives no title, identifier or complete bibliography entry. It describes the answer as well understood, calls rotational virtuals the natural domain, and reports Kauffman-bracket/Jones extensions to all virtuals; these are supplied assessments, not independently checked conclusions.
\item Registry status partially\_solved and classification PARTIAL-PROGRESS. Source evidence NEEDS\_REVIEW presents an open problem and calls for a current primary-literature audit.
\item Dataset attribution: Ulam AI / UnsolvedMath Contributors, CC BY 4.0. Source rights NEEDS\_REVIEW: publicly accessible author/publisher copy, redistribution terms requiring review.
\end{itemize}
\end{document}
